\section{Conclusions}
We formulate mapping-assisted semantic alignment for RF receiver reports. Known field meanings support physical comparison, and mapped class probabilities support event association. Receiver exclusivity and physical admission restrict the output partition, with experiments evaluating these feasibility requirements separately from event accuracy.

The mapping experiments hide three receiver maps while retaining a known reference to fix common class names. Independent adaptation scenes recover missing correspondences within the four-class vocabulary. Exact conditional search avoids a coordinate-update failure, while spectral synchronization provides comparable accuracy in the tested conditions. Alternating updates do not improve on sequential estimation. Wrong relative mappings harm association, whereas a shared class renaming leaves it unchanged.

Dense same-class collisions remain difficult after complete map recovery because correct column meanings cannot distinguish every transmission. Shared wrong predictions can also preserve event compatibility, separating semantic registration from class correctness and event identity. Integer programs quantify greedy objective loss and report-processing cost. LoED remains an auxiliary metadata check. Its timestamp audit exposes packet-reference ambiguity. A training-metadata clock estimate removes one use of reference identities in that check. Field validation still needs independent event truth and measured class probabilities from several receivers.
